\documentclass[11pt]{article}
\usepackage[T1]{fontenc}
\usepackage{lmodern}
\usepackage[margin=1in]{geometry}
\usepackage{amsmath,amssymb,amsthm,booktabs}
\usepackage[hidelinks]{hyperref}
\hypersetup{pdftitle={A better first step can delay stopping in a spectral gradient method},pdfauthor={},pdfsubject={AI-assisted working paper, pending author review}}
\newtheorem{theorem}{Theorem}
\newcommand{\R}{\mathbb R}
\newcommand{\norm}[1]{\lVert#1\rVert}
\newcommand{\bb}{\mathrm{ABBmin1}}
\title{A better first step can delay stopping in a spectral gradient method}
\author{}
\date{November 2026}
\begin{document}
\maketitle
\begin{abstract}
A spectral gradient method can retain an old secant after the associated
eigenmode has vanished from the current gradient. We give an exact example
for an ABBmin1 rule with a ten-secant window. An interval of alternative
first steps makes the gradient at the next point smaller than the
reference gradient. At the prescribed tolerance equal to half the common
post-initialization gradient norm, the reference stops at point five
and every alternative in the interval remains above tolerance through point five.
Three closed-form first steps lie within the interval. Two of them
stop earlier at a wider tolerance on the same objective.
The proof uses the original delayed secant to identify the precise
step-selection decisions. The result concerns gradient requests at a
fixed tolerance, with physical timing left open.
\end{abstract}

\noindent\emph{Research note.} This AI-assisted write-up reconstructs a
previously archived exact analysis \cite{source} and remains pending author review.

\section{Step rule and the experiment}

Let \(g_k=\nabla f(x_k)\) and \(x_{k+1}=x_k-\alpha_k g_k\).
Once \(g_k\) has been obtained, we test \(\norm{g_k}\le\epsilon\) before
selecting the next step. For a policy \(P\), its stopping index \(\tau_P(\epsilon)\) is the
first \(k\ge0\) that satisfies this test, with value \(+\infty\) if
the set is empty.
For a nonzero secant \(s_{k-1}=x_k-x_{k-1}\) and the corresponding
gradient difference \(y_{k-1}=g_k-g_{k-1}\), define
\[
 a^1_k=\frac{s_{k-1}^{\top}s_{k-1}}{s_{k-1}^{\top}y_{k-1}},
 \qquad
 a^2_k=\frac{s_{k-1}^{\top}y_{k-1}}{y_{k-1}^{\top}y_{k-1}}.
\]
Our reference uses the ABBmin1 family of Frassoldati, Zanni, and Zanghirati
\cite{fzz} with the following fixed parameters. At \(k\ge1\), it saves
\(a^2_k\). If \(a^2_k/a^1_k<4/5\), it takes the smallest \(a^2_j\)
for \(\max(1,k-9)\le j\le k\), with ties assigned to the earliest index.
Otherwise it takes \(a^1_k\). Write \(R\) for this reference run.
An alternative run \(P_\alpha\) performs the same initialization and
saves the same secant at \(k=1\). It takes the prescribed first step
\(\alpha_1=\alpha\), and from \(k=2\) onward it follows the reference
rule using its own stored secants. The analysis counts the gradients
that the stated rule actually requests. Related modifications act
later in a line search. One uses a trial function value before the
next gradient request~\cite{trial}. Another uses two endpoint
gradients to certify a stopping point inside the segment~\cite{chord}.
Here the only intervention is a prescribed first step after
initialization.

We use \(L=5000\), \(c=101/100\), and the notation
\[
 \begin{gathered}
 z=(1/\sqrt{200},\,1/\sqrt{600},\,\sqrt{(L-1)/4})^\top,\quad
 H=\operatorname{diag}(c,L,L)-zz^\top,\\
 u=(1,0,0)^\top,\qquad e=(-1/2,\sqrt3/2,0)^\top,\qquad G=(L-c)/2 .
 \end{gathered}
\]
On \(\R^3\times\R\) let
\[
 f(x,t)=\tfrac12 x^\top Hx+Gu^\top x+t^2+\tfrac1{100}(1-\cos t).
\]
Set \((x_0,t_0)=(-e,0)\) and use the common step \(\alpha_0=1/L\).
The fourth coordinate is stationary on the entire visited trajectory.

\begin{theorem}\label{thm:main}
At the tolerance \(\epsilon=G/2=499899/400\), the reference first
stops at point \(5\), whereas
\[
 \tau_{P_\alpha}(G/2)>5\qquad
       \text{for every }\alpha\in(3/25,7/50).
\]
For every such \(\alpha\), the norm at point \(2\) is strictly
smaller for \(P_\alpha\) than for the reference.
\end{theorem}

\section{The vanished mode and the saved secant}

Write \(D=\operatorname{diag}(1/100,L-1,L-1)\).
The identity \(z^\top D^{-1}z=3/4+1/[600(L-1)]<1\) gives
\(H-I=D-zz^\top\succ0\) by weighted Cauchy--Schwarz.
Also \(LI-H=\operatorname{diag}(L-c,0,0)+zz^\top\succeq0\).
The fourth curvature lies in \([199/100,201/100]\), and
the Hessian is globally \(1/100\)-Lipschitz in operator norm.

Since \(z^\top e=0\) and \(He=Le+Gu\), the initial gradient is
\(g_0=(-Le,0)\). The common initialization reaches the origin,
where \(g_1=(Gu,0)\). The first secant is therefore
\(s_0=(e,0)\), \(y_0=(Le+Gu,0)\), and its two steps are
\[
 a^1_1=\frac4{3L+c},\qquad
 \delta:=a^2_1=\frac{3L+c}{3L^2+c^2}.
\]
Here \(1/5000<\delta<2001/10^7\) and \(a^2_1/a^1_1<4/5\).
For the last inequality, cross-multiplication gives
\(3L^2-30cL+11c^2>0\), which follows from \(L>10c\).
The reference accordingly takes \(\delta\) at \(k=1\).

Let \(w=(0,z_2,z_3)^\top/\sqrt{z_2^2+z_3^2}\). From point 1 on,
the gradients lie in the invariant plane with orthonormal basis
\((u,w)\), where \(H\) has matrix
\[
 K=\begin{pmatrix}q&-b_K\\-b_K&d\end{pmatrix},\qquad
 q=\frac{201}{200},\quad b_K^2=\frac{749851}{120000},\quad
 d=\frac{2250149}{600}.
\]
A perpendicular direction \((0,z_3,-z_2)^\top\) is an
\(L\)-eigenvector. The first secant has nonzero projection on
this direction, even though every gradient from \(g_1\) onward
has zero projection there. Let \(\lambda_-<\lambda_+\) be the
eigenvalues of \(K\). Because \(H-I\succ0\) and \(b_K\ne0\),
the Rayleigh quotient at \(u\) gives \(1<\lambda_-<q<101/100\).
The Rayleigh quotient at \(w\) gives \(\lambda_+>d>3750\),
and the trace then gives \(\lambda_+<3751\).

Every post-initialization secant has an \(a^2\) value at least
\(1/\lambda_+\). Indeed, \((s^\top Hs)/(s^\top H^2s)\) is a
weighted mean of the inverse eigenvalues on the plane.
Since \(1/3751>\delta\), the saved \(a^2_1\) is the unique
minimum whenever it belongs to the window. In particular, this
holds at \(k=2,3,4\) for both the reference and each alternative.

\section{Exact decisions and the stopping order}

Let \(U_k,V_k\) be the two components of \(g_k/G\) in the
low and high eigendirections, with signs selected so that
\(U_1,V_1>0\). Let \(Z_k=(V_k/U_k)^2\). The high-mode squared
projection at \(u\) is
\[
 p=V_1^2=\frac{b_K^2}{(\lambda_+-q)^2+b_K^2}.
\]
Using \(156/25<b_K^2<25/4\) and \(3748<\lambda_+-q<3750\)
gives \(44/10^8<p<45/10^8\), hence
\[
 \frac{44}{10^8}<Z_1=\frac{p}{1-p}<\frac{46}{10^8},
 \quad U_1>\frac{999}{1000},
 \quad V_1<\frac1{1400}.
\]
For a step \(\delta\), set \(A=1-\delta\lambda_-\),
\(B=1-\delta\lambda_+\) and \(\chi=(B/A)^2\). The eigenvalue
and secant bounds give
\[
 \frac{9997}{10000}<A<1,\qquad \frac{249}{1000}<B<\frac14,\qquad
 \frac{31}{500}<\chi<\frac{63}{1000}.
 \tag{1}\label{eq:chi}
\]
The components are multiplied by \(A,B\), and \(Z\) by \(\chi\),
whenever the step actually selected is \(\delta\).

At \(k\ge2\), the secant starts with \(g_{k-1}\): it is
\(s_{k-1}=-\alpha_{k-1}g_{k-1}\), \(y_{k-1}=Hs_{k-1}\)
on the plane. Substitution gives
\[
 A_1(Z)=\frac{1+Z}{\lambda_-+\lambda_+Z},\qquad
 \mathcal Q(Z)=
 \frac{(\lambda_-+\lambda_+Z)^2}
      {(1+Z)(\lambda_-^2+\lambda_+^2Z)},\qquad Z=Z_{k-1}.
 \tag{2}\label{eq:ratios}
\]
The chosen step acts on \(g_k\). Retaining the index \(k-1\)
in \eqref{eq:ratios} is necessary for both trajectories below.

\paragraph{Reference decisions.}
At \(k=2\), the ratio is at most
\(1.012^2/(1+3750^2\cdot44/10^8)<4/5\), so the
window supplies \(\delta\). From \eqref{eq:chi},
\(27/10^9<Z_2<29/10^9\). At \(k=3\), the ratio is at most
\(1.011^2/(1+3750^2\cdot27/10^9)<4/5\), so the step is
again \(\delta\). Then \(Z_3<(63/1000)^2(46/10^8)<2/10^9\).
The ratio read at \(k=4\) is greater than
\(1/(1.001\cdot1.03)>4/5\), by keeping \(\lambda_-^2\)
in the numerator and using \(\lambda_+/\lambda_-<3751\)
in the denominator. Thus \(\alpha_4=A_1(Z_3)\).
Points \(1,\ldots,4\) all have \(U_k=U_1A^{k-1}>0.99\).
At point \(5\) the two components instead satisfy
\[
 |U_5|\le 3751\,Z_3\,|U_4|<10^{-5},\qquad
 |V_5|\le 3751\,|V_4|
       <\frac{3751}{1400\cdot64}<\frac1{20}.
\]
These bounds use \(1-A_1(Z_3)\lambda_-
=(\lambda_+-\lambda_-)Z_3/(\lambda_-+\lambda_+Z_3)\)
and \(1-A_1(Z_3)\lambda_+
=(\lambda_--\lambda_+)/(\lambda_-+\lambda_+Z_3)\).
Together with \(\norm{g_0}=L\), they give the first hit at point \(5\).

\paragraph{Alternative decisions.}
For \(\alpha\in(3/25,7/50)\),
\[
 17/20<1-\alpha\lambda_-<9/10,\qquad
 449<\alpha\lambda_+-1<525 .
\]
Thus \(Z_2=Z_1[(\alpha\lambda_+-1)/(1-\alpha\lambda_-)]^2\)
lies in \((1/10,1/5)\). At \(k=2\) the ratio still reads the
common \(Z_1\) and the window again supplies \(\delta\).
Put \(t=\lambda_-/\lambda_+<1/1000\). The bound
\(\mathcal Q(Z)=(Z+t)^2/[(1+Z)(Z+t^2)]\le (Z+t)^2/Z\)
gives \(\mathcal Q(Z_2)<0.201^2/0.1<4/5\).
The step at \(k=3\) is consequently \(\delta\) as well.
Since \(0.006<Z_3=\chi Z_2<0.013\), the ratio at \(k=4\) is
less than \(0.014^2/0.006<4/5\), which gives the third
consecutive \(\delta\) after the override. Therefore
\[
 U_k=U_1(1-\alpha\lambda_-)A^{k-2}
     >0.999\cdot0.85\cdot0.9997^3>0.84
       \quad (2\le k\le5).
\]
The common gradients at points \(0,1\) also exceed \(G/2\),
proving the strict comparison of stopping indices.

\paragraph{Comparison at point 2.}
For a first step \(\gamma\) from point 1,
\[
 \frac{\norm{g_2(\gamma)}^2}{G^2}
   =1-2q\gamma+(q^2+b_K^2)\gamma^2.
 \tag{3}\label{eq:norm}
\]
The quadratic decreases on \([0,q/(q^2+b_K^2)]\).
For the interval in the theorem, compare its two values directly.
It is \((\alpha-\delta)\{(q^2+b_K^2)(\alpha+\delta)-2q\}\).
Since \(q^2+b_K^2<8\), \(\alpha+\delta<0.141\) and \(2q>2\),
this difference is negative. This completes the proof of
Theorem~\ref{thm:main}.

\section{Three closed-form first steps and two tolerances}

We next locate three previously constructed curvature-based first
steps in the certified interval. This comparison keeps all second
and later actions as defined above. Use the valid Hessian Lipschitz bound \(M=1/24\), which is
looser than \(1/100\) for this instance, and set
\(h=1/\sqrt{1200}\), \(H_0=1+h\). Define
\[
 C_A(r)=\int_0^r\min\{L,A+Mb\}\,db,\qquad
 \mathcal R_A=C_A^{-1},\qquad \rho=\mathcal R_{H_0}(G).
\]
These definitions come from the curvature envelope along the current
negative-gradient ray after the common initialization step \(1/L\).
The angle cosine between \(e\) and \(-u\) is \(1/2\), the old chord length is \(1\), and
\(1/100\) is an upper bound on the chord-average excess curvature
in direction \(u\). The active
conditions \(1/100<M(1/2)(1)/2\) and \(H_0+M\rho/2<L\) hold.
Indeed, \(1<H_0<103/100\), the cap is inactive for \(r\le330\),
and \(C_{H_0}(300)<2184<G<10395/4<C_{H_0}(330)\).
Thus \(300<\rho<330\). On the uncapped branch, solving
\(Ar+Mr^2/2=z\) gives
\(\mathcal R_A(z)=2z/(A+\sqrt{A^2+2Mz})\).

For \(b\ge0\), let
\[
 \sigma(b)=\sqrt{1+b^2+b}-b,\qquad
 \omega(b)=h(1-\sqrt{\sigma(b)}),\qquad \beta=\rho/4.
\]
Here \(0<\omega(b)<h\) for \(b>0\) and \(\omega\) is
strictly increasing. To see this, the derivative of
\(\sigma\) is \((2b+1)/(2\sqrt{1+b^2+b})-1<0\).

For completeness, let \(q_f(b)\) be the curvature along \(u\) at
the future point \((-bu,0)\). On the old chord \((-te,0)\),
\(0\le t\le1\), the average curvature along \(u\) is \(q\)
in this objective, so its excess over 1 is at most \(a=1/100\).
The same upper bound follows using the two gradients alone.
The chord Hessian average \(\overline H\) satisfies
\((LI-\overline H)e=-Gu\) after the common \(1/L\) step.
Cauchy--Schwarz for \(LI-\overline H\succeq0\) then gives
\(u^\top\overline H u\le L-G^2/(-e^\top Gu)=c\).
The distance \(\norm{te-bu}=\sqrt{t^2+b^2+tb}\) is at most
\(b+\sigma(b)t\), its endpoint-interpolation bound on \([0,1]\).
The Hessian Lipschitz inequality makes the old excess curvature
at least \(\max\{0,q_f(b)-1-Mb-M\sigma(b)t\}\).
If \(q_f(b)-1-Mb>h\sqrt{\sigma(b)}\), its integral is strictly
larger than \(a\): at equality it is a triangle of area
\(h^2/(2M)=a\), with support length
\(\sqrt{2a/(M\sigma(b))}\le1\) because \(\sigma(b)>1/2\).
Thus \(q_f(b)\le\min\{L,1+Mb+h\sqrt{\sigma(b)}\}\).
Where the base envelope is uncapped, this improves it by
\(\omega(b)\). This explains the three scalar constructions.
Put \(v=\beta\omega(\beta)\),
\(w=(\rho/8)\omega(\rho/8)\),
\(z_d=(\rho/16)\omega(\rho/16)\), and \(d_E=\omega(\beta)\).
The three radii and the corresponding steps are
\[
 \begin{aligned}
 r_B&=\mathcal R_{H_0}(G+v+w),\\
 r_C&=\mathcal R_{H_0}(G+v+w+z_d),\\
 r_E&=\mathcal R_{H_0-d_E}(G-d_E\beta),\qquad
 \alpha_P=r_P/G\quad(P=B,C,E).
 \end{aligned}
\]
The quantities \(v,w,z_d\) integrate constant lower bounds of the
curvature deficit over \([\rho/4,\rho/2]\),
\([\rho/8,\rho/4]\), and \([\rho/16,\rho/8]\), respectively.
These intervals only meet at endpoints. The rule for \(r_E\) uses a
piecewise upper envelope: its intercept is lowered by \(d_E\)
from \(b=\beta\) onward. The old and lowered slopes are the same
in the uncapped region. Consequently for \(r\in[\beta,330]\)
its integral is \(C_{H_0-d_E}(r)+d_E\beta\).

Since \(v+w+z_d<7h\rho/16<5\), we have
\(\rho<r_B<r_C<330\). The integral of the single-breakpoint
envelope at \(r_C\) is
\(C_{H_0}(r_C)-d_E(r_C-\beta)<G\), because
\(v+w+z_d<(7\rho/16)d_E\) and \(r_C-\beta>3\rho/4\).
At \(330\) the same integral is strictly greater than \(10395/4\)
as \(H_0-d_E>1\). Hence \(r_C<r_E<330\). Since \(2499<G<2500\),
all three actions belong to \((3/25,7/50)\).

For the wider tolerance, let
\(n_P=\norm{g_2(\alpha_P)}\) and \(\epsilon_0=(n_C+n_B)/2\).
The squared gradient norm as a function of the radius \(r=G\gamma\)
decreases until \(Gq/(q^2+b_K^2)>330\), by direct rational
comparison of the specified \(G,q,b_K^2\).
As \(0<G\delta<1<r_B<r_C<r_E<330\), we have
\[
 n_E<n_C<\epsilon_0<n_B<\norm{g_2(\delta)}<G<L .
\]
Thus \(C\) and \(E\) first stop at point \(2\) on this fixed
wider tolerance, and \(R,B\) have not stopped there.
The tolerance is defined by this fixed instance. In contrast, in
Theorem~\ref{thm:main} the tolerance \(G/2\) is fixed from
the common gradient before any of the branches. The two
tolerances yield opposite orders on the same objective.

\section{Costs and scope}

The reference obtains six gradients and performs five actual
updates to meet \(G/2\). The other trajectories have spent the
same six gradient requests by point \(5\) without meeting it.
If they finish at a finite index, they need at least one more request.
Every branch evaluates its own norms and secant ratios. The
additional work for a curvature-based first step includes supplying
its effective curvature bounds, computing its gate and scalar
root, and retaining the secant that the step overrides.
The eigenvalues are used to analyze the iteration. The
algorithm requests the full gradient and computes its steps
from the stored secants.

The proof uses exact arithmetic. The strict inequalities identify
where a numerical implementation would need error control.
The theorem characterizes a finite stopping-order reversal
caused by a retained secant. Wall-clock cost, the eventual
behavior of an alternative run, and the corresponding behavior
on nonquadratic trajectories remain open.

\begin{thebibliography}{9}
\bibitem{fzz}
G. Frassoldati, L. Zanni, and G. Zanghirati.
New adaptive stepsize selections in gradient methods.
\emph{Journal of Industrial and Management Optimization}, 4(2):299--312, 2008.
\href{https://doi.org/10.3934/jimo.2008.4.299}{DOI: 10.3934/jimo.2008.4.299}.
\bibitem{source}
\emph{Stopping Targets and Cold-Window Memory} (in Chinese).
Private research report, October 24, 2026. Exact benchmark
and window mechanism. The archive preserves the source and its
independent proof records.
\bibitem{trial}
\emph{Reselecting steps from line-search trial values} (in Chinese).
Working paper, November 2026.
\href{https://taitai2025.github.io/assets/documents/research/trial-firsthit.pdf}{PDF}.
\bibitem{chord}
\emph{A certified stopping point from two gradient samples}.
Working paper, November 2026.
\href{https://taitai2025.github.io/assets/documents/research/certified-chord.pdf}{PDF}.
\end{thebibliography}
\end{document}

